% Section 4: DoF and Synthesis Theorem v3
\section{Degrees of Freedom and the Synthesis Theorem}
\label{sec:synthesis-main}

\subsection{Motivation}

The unified language provides a common representation for structure,
adversarial behavior, and dynamics. Equivalence does not hold universally.
The DoF invariant characterizes precisely when it emerges.

\subsection{Degrees of freedom}
\label{subsec:dof-main}

\paragraph{Representation DoF.}
$\mathrm{DoF}_{\mathrm{rep}}(P)=0 \iff \nexists\,P'\neq P:F(P')\subseteq P'.$

\paragraph{Source DoF.}
$\mathrm{DoF}_{\mathrm{src}}(P):=1/C_{\min}(P).$ Zero iff $C_{\min}=+\infty$.

\paragraph{Dynamic DoF.}
$\mathrm{DoF}_{\mathrm{dyn}}(P)=0 \iff P$ is the unique stable attractor of $\mathcal{A}$.

\begin{definition}[DoF vector and Zero-DoF Condition]\label{def:dof-main}
$\mathrm{DoF}(P):=(\mathrm{DoF}_{\mathrm{rep}},\mathrm{DoF}_{\mathrm{src}},\mathrm{DoF}_{\mathrm{dyn}})$.
Zero-DoF: $\mathrm{DoF}(P)=(0,0,0)$.
\end{definition}

\subsection{The Synthesis Theorem}
\label{subsec:synthesis}

The three notions --- correctness, unbreakability, stability --- are projections
of the same structure onto three axes of $(F,\mathcal{B},\mathcal{A})$.
The shared semantic object is $\mathrm{SC}_T$: what all three disciplines
call canonicality, in different terms.

\begin{theorem}[CRT--IET Synthesis]\label{thm:synthesis-main}
Assume C1--C6 ($\mathrm{UES}(P)=\emptyset$); IET Paper~1~v2.2 conditions;
and shared $\mathrm{SC}_T$. Then:
$$\boxed{\mathrm{DoF}(P)=(0,0,0)
  \;\Longleftrightarrow\;
  \mathrm{SC}_T(P)=1
  \;\Longleftrightarrow\;
  C_{\min}(P)=+\infty
  \;\Longleftrightarrow\;
  \mathrm{Stable}_{\mathcal{A}}(P)=1.}$$
\end{theorem}

\begin{proof}
\textbf{DoF translation (primary).}
$\mathrm{DoF}_{\mathrm{rep}}=0\iff\mathrm{SC}_T=1$;
$\mathrm{DoF}_{\mathrm{src}}=0\iff C_{\min}=+\infty$;
$\mathrm{DoF}_{\mathrm{dyn}}=0\iff\mathrm{Stable}=1$.

\textbf{CRT bridge} (\cite{bridge-hierarchy}, C1--C6):
$\mathrm{SC}_T=1\iff C_{\min}=+\infty$.

\textbf{IET bridge} (\cite{iet-paper1}):
$\mathrm{SC}_T=1\iff\mathrm{Stable}=1$.

Combining: all four quantities are equivalent to $\mathrm{DoF}(P)=(0,0,0)$.\qed
\end{proof}

\begin{remark}[DoF collapse, not coincidence]
The equivalence is a structural collapse: when no degree of freedom remains,
there is only one thing the system can be --- and all three disciplines
arrive at the same answer.
\end{remark}

\begin{remark}[SC as alignment point, not chain]
$\mathrm{SC}_T$ is the shared semantic object. The proof is not a chain
$C_{\min}\to\mathrm{SC}_T\to\mathrm{Stable}$; both bridges meet at $\mathrm{SC}_T$.
\end{remark}

\subsection{Failure and boundary}\label{subsec:boundary}

$\mathrm{UES}(P)=\emptyset\iff C_{\min}(P)=+\infty\iff\mathrm{DoF}_{\mathrm{src}}(P)=0$.

\begin{center}\begin{tabular}{llp{4cm}}\hline
$\mathrm{DoF}(P)$ & System & Result\\\hline
$(0,0,0)$ & Class I/II & Full equivalence\\
$(0,{>}0,0)$ & Class III (DeFi) & $C_{\min}\geq\theta$ only\\
$(0,{>}0,{>}0)$ & Open multi-eq & Both fractured\\\hline
\end{tabular}\end{center}

\begin{quote}\itshape
The synthesis theorem is not an accidental equivalence.
It is a collapse induced by the elimination of all degrees of freedom.
\end{quote}
